\section{Related Work}
\subsection{Truy vấn liên tục và tính hợp lý của đầu ra}
TransProtect\cite{transprotect} khai thác dự đoán chuyển động và chi phí tới đích để chọn miền obfuscation,
đối phó tấn công biết bối cảnh đường/giao thông. Semantic correlation\cite{semantic} xét sự nhất quán
ngữ nghĩa và thời gian của chuỗi tập vị trí. Fake queries\cite{fake} chèn thêm tập toàn giả giữa các lần
truy vấn để làm gián đoạn tương quan và reachability. Những hướng này chỉ ra vì sao attacker không nên
bị giới hạn vào một điểm độc lập; chúng khác nhau ở việc có công bố GPS thật như một thành viên hay không.

\subsection{Geo-I, utility và các phạm vi bảo vệ khác}
EV querying\cite{ev} kết hợp approximate Geo-I với dummy và xử lý cục bộ để giữ chất lượng tìm trạm sạc;
điều đó cho thấy Geo-I + dummy + local selection đã có tiền lệ. PTPPM\cite{ptppm} là preprint road-aware
kết hợp personalized budget và distortion privacy. Khác biệt cần phân tích của GeoI-BQ nằm ở đặc tả
kênh Q và cận toàn transcript, không ở phát biểu rằng lần đầu dùng nhiễu trên mạng đường.

CPCROK\cite{cpcrok} bảo vệ việc nối lại beacon qua pseudonym trong VANET; DP-FETC\cite{dpfetc} phát hành
dataset tổng hợp và xử lý OD/tương quan. Các mục tiêu đó liên quan identity/endpoint nhưng không tự
tương đương tác vụ POI online. SoK\cite{sok} giúp phân biệt đơn vị privacy và các mục tiêu của generative
trajectory publication; nó là nguồn bối cảnh, không là một model đối chứng online.

\subsection{Tổng hợp phạm vi, metrics và khả năng đối chiếu}
Khảo sát phương pháp gần đây dùng cửa sổ \textbf{28/10/2023--28/10/2026}, xét công bố đầu tiên có bằng chứng.
Geo-I 2013 và predictive mechanism 2014 được dẫn riêng ở Phần 3 như nguồn nền tảng.
Trong bảng: \textbf{F} = trực tiếp xét/đánh giá task theo giả định paper; \textbf{P} = một phần hoặc bối cảnh khác;
\textbf{CX} = chưa đủ căn cứ đối chiếu chi tiết. Các nhãn do ta ánh xạ sang S1--S10, không có nghĩa chống mọi attacker.
Scenario không liệt kê chưa có căn cứ coverage trong audit này; không tự quy thành score 0 hoặc đã bảo vệ.

\begin{footnotesize}\par\smallskip\noindent\begin{tabularx}{\linewidth}{@{}P{2.75cm}P{2.15cm}P{1.6cm}P{2.4cm}P{2.5cm}Y@{}}\toprule
\textbf{Paper / cơ chế} & \textbf{Output} & \textbf{Scenario} & \textbf{Privacy} & \textbf{Utility / cost} & \textbf{Tái lập trong bản này} \\
\midrule
\textbf{TransProtect} 2024\cite{transprotect}\newline Học chuyển động + chọn miền + obfuscation & Một tọa độ thay thế/mốc & F: S1, S3\newline P: S2 & EIE; tấn công VehiTrack & Travel-cost distortion $\Delta$c; chi phí vận hành chưa có contract chung & Chạy Markov adaptation; không có checkpoint Transformer gốc \\
\textbf{Semantic correlation} 2026\cite{semantic}\newline Tương quan ngữ nghĩa và thời gian & GPS thật + K-1 dummy & F: S1\newline P: S2, S3 & Anonymity success rate (ASR) & DER là chất lượng dummy; thời gian sinh & Chạy OSM/empirical adaptation; không tái lập AMap/LSTM \\
\textbf{Fake queries} 2026\cite{fake}\newline Chèn query toàn giả giữa query thật & Real+dummy sets và fake packets & F: S1, S3\newline P: S2 & ASR; số đường khó phân biệt & Thời gian/RSS; chưa đo Recall POI chung & Chạy road/history adaptation, giữ fake timing và mọi chi phí \\
\textbf{EV querying} 2024\cite{ev}\newline Approximate Geo-I + dummy + Edge & Vị trí nhiễu và dummy; chọn trạm sạc local & P: S1, S2, S3 & ($\varepsilon$,$\delta$)-AGeoI; IBU/Wasserstein & Cost of Privacy (CoP), tỷ lệ CoP=0 & Khảo sát; chưa có adaptation cùng dịch vụ và giả định Edge \\
\textbf{CPCROK} 2025\cite{cpcrok}\newline Mix-zone, đổi pseudonym và beacon giả & Beacon thật + tuyến beacon giả & P: S4 & Tỷ lệ khớp lại quỹ đạo $\rho$ & Số tuyến giả $\psi$; truyền thông VANET & Khác hợp đồng POI online; không đưa vào bảng score chính \\
\textbf{DP-FETC} 2025\cite{dpfetc}\newline Thống kê DP, tổng hợp trace và chỉnh OD & Bộ dữ liệu quỹ đạo offline & P: S3, S8, S9, S10 & $\varepsilon$-DP; mutual information & JSD phân bố; độ phức tạp & Offline publishing; không chấm bằng request POI online \\
\textbf{PTPPM} 2025\cite{ptppm}\newline Personalized budget, PLS, Permute-and-Flip & Vị trí đã perturb trên mạng đường & CX: S1, S3 & Geo-I + distortion privacy/EIE bound & Perturbation distance/QoS & Preprint liên quan; chưa có comparator đã audit trong repo \\
\textbf{SoK trajectory generation} 2024\cite{sok}\newline Systematization và unit of privacy & Khảo sát phát hành dữ liệu tổng hợp & Không áp dụng & Formal privacy, reconstruction/correlation threats & Mục tiêu và utility của dữ liệu tổng hợp & Nguồn bối cảnh; không là một model đối chứng online \\
\bottomrule\end{tabularx}\par\smallskip
\end{footnotesize}

\textbf{Khoảng trống đánh giá.} ASR đòi một thành viên thật và posterior tương ứng; Q-only không có thành viên
đó. DER đo chất lượng dummy, JSD đo fidelity của dataset và CoP đo tổn thất dịch vụ, không phải cùng một
đại lượng privacy. Vì vậy, Phần 4 định nghĩa metrics chung theo threat và utility thực nhận, đồng thời giữ
metrics gốc đúng contract. Khác định nghĩa phải ghi rõ; thiếu đầu vào giữ N/A, không chuyển thành score tốt.

Ba hàng TransProtect*, Semantic* và Fake-query* trong thực nghiệm là các adaptation đã audit. Những hàng
khác định vị nghiên cứu và giới hạn coverage; chưa có implementation phù hợp thì không được gán số benchmark.
